\section{Un langage de représentation -- XML}

\begin{frame}
\frametitle{XML, format pour données ``semi-structurées''}

XML est d'abord un standard du World-Wide-Web Consortium (W3C).

\begin{itemize}
  \item La sérialisation de documents XML est normalisée, ce qui permet
       des échanges réseau sans perte d'information.
 
  \item XML est un format générique, qui se spécialise ensuite en
     \alert{``dialectes''} pour des domaines spécifiques (e.g., XHTML,
     affichage dans un navigateur).

  \item Le W3C a défini des standards associés: DOM (le ``modèle''), 
   XSchema (le typage), XPath (navigation), XSLT (restructuration),
     XQuery (requêtes), et beaucoup d'autres. Très compliqué, et surtout
     très lourd.
\end{itemize}

\begin{remark}
XML est une version simplifiée de \alert{SGML}.
\end{remark}

\end{frame}



\begin{frame}[fragile]
\frametitle{Documents XML: structure et contenu}

XML permet de \alert{structurer} du contenu textuel.

\vfill

Le contenu textuel:

\begin{center}
  Le film de SF Gravity,  
  réalisé par Alfonso Cuaròn, est paru en l'an 2013.
\end{center}

\vfill

Le même contenu, structuré en XML.

\begin{lstlisting}
    <?xml version="1.0" encoding="utf-8"?>

    <film>
      <titre>Gravity</titre>
      <annee>2013</annee>
      <genre>SF</genre>
      <realisateur
         nom="Cuaron"
         prenom="Alfonso"/>
    </film>
\end{lstlisting}

\end{frame}


\begin{frame}[fragile]
\frametitle{Documents XML: modèle et sérialisation}

\begin{description}

\item [Modèle] Un document XML est un arbre, chaque n{\oe}ud a un type particulier.

\item [Sérialisation] tout document XML peut être transformé en une chaîne
   de caractères représentant l'arborescence du document.

\end{description}

La norme XML définit une \alert{syntaxe} pour coder des documents textuels
hiérarchiques. Aucune interprétation n'est
définie à priori.

\vfill

Les dialectes XML imposent des contraintes, et définissent une interprétation.

\vfill

\begin{beamerboxesrounded}{Le dialecte XHTML}
Un document XHTML est un document XML. Il obéit à des contraintes de structure, et s'interprète
sous forme de règles d'affichage à l'écran.
\end{beamerboxesrounded}

\end{frame}

\begin{frame}[fragile]
\frametitle{Qui utilise XML?}

XML est le standard (W3C) pour la représentation  de données textuelles.

\vfill
\alert{Exemples}\,: XHTML, DocBook, RSS, ...

\vfill

XML est AUSSI utilisé pour sérialiser des données applicatives, par exemple pour les données
gérées par vos applications personnelles.


\vfill

\alert{Exemples}\,: Traitements de texte (Word), tableurs (Excel), calendriers, graphiques (en SVG), etc.

\vfill

\begin{beamerboxesrounded}{Dans tous les cas}
La représentation en XML permet de \alert{réutiliser} le contenu, de \alert{l'échanger} avec d'autres applications.
\end{beamerboxesrounded}

\end{frame}

\begin{frame}[fragile]{Exemple: flux de données RSS}

Pour disséminer des informations en mode "push" (cf. Blogs, sites de news, etc.).

\begin{lstlisting}[language=myXML]
<?xml version="1.0" encoding="UTF-8" ?>

<rss version="2.0">
<channel>
	<title>Webdam Project</title>
	<atom:link href="http://webdam.inria.fr/wordpress/?feed=rss2" 
	            rel="self" type="application/rss+xml" />
	<link>http://webdam.inria.fr/wordpress</link>
	<description>Web Data Management</description>
	<pubDate>Wed, 26 May 2010 09:30:54 +0000</pubDate>

  <item>
    <title>News for the beginning of the year</title>
    <description>...</description>
    <link>http://webdam.inria.fr/wordpress/?p=475</link>
    <pubDate>Fri, 15 Jan 2010 08:48:45 +0000</pubDate>
    <dc:creator>Serge</dc:creator>
    <category>News</category>
  </item>
</channel>
</rss>
\end{lstlisting}
\end{frame}


\begin{frame}[fragile]
\frametitle{Scalable Vector Graphics (SVG)}

Et oui, on peut représenter des données graphiques en XML !

\bigskip

\begin{tabular}{m{.7\hsize}|m{.2\hsize}}
\begin{small}

\begin{lstlisting}
<?xml version="1.0" encoding="UTF-8" ?>
<svg xmlns="http://www.w3.org/2000/svg">

   <polygon points="0,0  50,0  25,50"
        style="stroke:#660000;"/>

  <text x="20" y="40">Some SVG text</text>
</svg>
\end{lstlisting} 
\end{small}
&
\includegraphics{svg}
\end{tabular}
\end{frame}


\begin{frame}[fragile]

\frametitle{Music XML}

Ou des partitions musicales.

\medskip

\begin{tabular}{m{.6\hsize}|m{.3\hsize}}
\begin{small}

\begin{lstlisting}
<?xml version="1.0" encoding="UTF-8"?>

<score-partwise version="2.0">
  <part id="P1">
      <attributes>
        <divisions>1</divisions>
      </attributes>
      <note>
        <pitch>
          <step>C</step>
          <octave>4</octave>
        </pitch>
        <duration>4</duration>
      </note>
  </part>
</score-partwise>
\end{lstlisting}
\end{small}

&
\includegraphics{musicxml}

\end{tabular}
\end{frame}



\section{Formalisme XML: l'essentiel}



\begin{frame}[fragile]
\frametitle{Syntaxe XML}

Un document XML est composé de :

\vfill

\alert{Un prologue}

\begin{enumerate}
\item déclaration de document XML \verb7<?xml version = "1.x" [encoding = "norme" ]? >7 

\item suivie éventuellement (sans contrainte d’ordre) 

\item d’une référence à une DTD \verb7<!DOCTYPE element racine SYSTEM "uri" >7 
\item de commentaires \verb?<!-- texte --> ?
\item d’instructions de traitement (par ex. appel à XSLT) \verb7<?instruction ... ? >7 

\end{enumerate}

\vfill

\alert{Un corps}

\begin{itemize}
\item une imbrication d'\alert{éléments} et de \alert{textes}
\end{itemize}

\end{frame}


%%%%%%%%%%%%%%%%%%%%%%%%%% BEGIN PHILIPPE



\begin{frame}[fragile]
\frametitle{Quelques exemples de base}

Remarque: les retours à la ligne n'ont aucune signification.
\begin{smallverbatim}
  <document/>
  
  <document> Hello World! </document>
  
  <document>
    <salutation> Hello World! </salutation> 
  </document>
  
  <?xml version="1.0" encoding="utf-8" ?> 
  <document>
    <salutation color="blue"> Hello World! </salutation> 
  </document>
\end{smallverbatim}

Le dernier exemple est le seul complètement correct\,: il comprend le prologue.

\end{frame}


\begin{frame}[fragile]
\frametitle{Forme hiérarchique des documents\,: le modèle DOM}

Un document XML est interprété comme un \alert{arbre}, basé sur
 un modèle, le DOM (\alert{Document Object Model}) normalisé par le W3C.

\vfill

Dans le DOM, chaque n{\oe}ud de l'arbre a un \alert{type}, et une \alert{description} (nom, contenu)

\vfill

Les n{\oe}uds de type \domtype{Element} correspondent au balisage dans la forme sérialisée.

\begin{enumerate}
  \item ils définissent la \alert{structure} du document.
  \item ils ont un \alert{nom} (mais pas de \alert{valeur})
  \item ils ont un \alert{contenu}\,: le sous-arbre dont ils sont racine.
\end{enumerate}
\vfill
Les n{\oe}uds de type \domtype{Text} correspondent au \alert{contenu textuel} du document.

\begin{enumerate}
  \item ils constituent les feuilles de l'arborescence (un n{\oe}ud \domtype{Text} n'a pas de fils);
  \item ils ont une \alert{valeur} (mais pas de \alert{nom}).
  \end{enumerate}
\end{frame}

\begin{frame}[fragile]
\frametitle{Correspondance forme sérialisée - forme DOM}


Voici deux exemples de base illustrant la construction élément/texte.

\begin{tabular}{m{.6\hsize}m{.4\hsize}}
    
\begin{tinyverbatim}
<monElt>
  Contenu textuel.
</monElt>
\end{tinyverbatim}

 & 
\begin{tinyverbatim}
<maRacine>
  Contenu textuel
  <sousEl>
    Autre contenu
  </sousEl>
</maRacine>
\end{tinyverbatim}
\end{tabular}
 
 \vfill
 
Et la représentation arborescente (DOM).
\smallskip
 
\includegraphics[width=10cm]{ElementEtText}
 
\end{frame}

\begin{frame}[fragile]
\frametitle{Quelques remarques en vrac (anecdotique)}

\begin{itemize}
\item Les sous-éléments sont ordonnés.
\item Pas de caract\`eres spéciaux (mais \verb?-?,\verb?_?,\verb?.?,\verb?/? autorisés) dans les noms des éléments, pas d'espaces.
\item Un élément peut \^etre vide: \verb?<nom_elt></nom_elt>? 

  Notation abrégée\,: \verb?<nom_elt/>? (Exemples d'éléments vides (XHTML): \verb?<br/>? ou \verb?<img width="100" height="57" src="logoCnam.jpg"/>?) 
\item (Contrairement à HTML,) XML est \textit{case sensitive} donc \verb?<genie>? $\neq$ \verb?<Genie>? 
\item Les caractères spéciaux peuvent \^etre remplacés par une \alert{référence}:
\verb?&? doit \^etre remplacé par sa référence \verb?&amp;? (version numérique\,: \verb?&#28;?) \\
\verb?<? devient \verb?&lt;? (ou \verb?&#60;?)\\
\verb?>? devient \verb?&gt;? (ou \verb?&#62;?)\\
\verb?"? devient \verb?&quot;? (ou \verb?&#39;?)\\
\verb?'? devient \verb?&apos;? (ou \verb?&#34;?)\\

\end{itemize}

\end{frame}


%%%%%%%%%%%%% 
\begin{frame}[fragile]
\frametitle{D'autres types de n{\oe}ud\,:  \textbf{Attribute} et \textbf{Document}}



\begin{tabular}{m{.4\hsize}m{.5\hsize}}

Les \alert{attributs} sont des paires clé/valeur attachées à un élément.

Un attribut a une \alert{valeur}\,: un chaîne de caractères.

 & 
 
 
Un document XML sous forme sérialisée commence \alert{toujours} par un prologue
représenté par un nœud de type \textbf{Document}.
\smallskip
\begin{center}
\xmlinline{<?xml version="1.0"  encoding="utf-8" ?>}
\end{center}

\end{tabular}

\vfill

\includegraphics[width=10cm]{GravityDOM}
 

\end{frame}

\begin{frame}
\frametitle{Résumé de l'essentiel}


\textbf{Forme sérialisée}
\begin{itemize}
  \item Un document débute par un prologue
  \item Le contenu est enclos dans une unique balise ouvrante/fermante.
  \item Chaque balise ouvrante   \texttt{<nom>}
      \alert{doit} avoir une balise fermante \texttt{</nom>}; tout ce qui est entre les deux
        est soit du texte, soit du balisage correctement ouvert/fermé.
\end{itemize}

\textbf{Forme arborescente (DOM)}
\begin{itemize}
  \item Un document est un \alert{arbre} avec une racine (du document) de type \nodetype{Document}
        
  \item La racine du document a un seul \alert{élément} fils, de type \domtype{Element}, appelé
      \alert{l'élément racine}.
             
  \item Chaque \domtype{Element} est un sous-arbre représentant du contenu structuré
     formé de nœuds \domtype{Element}, de nœuds \domtype{Text}, et parfois
     de nœuds \domtype{Attribute}.
\end{itemize}

\end{frame}



\begin{frame}{Document bien formé}

Un document est bien formé s'il respecte la syntaxe XML. 


\begin{itemize}
\item \`a toute balise ouvrante est associée une balise fermante,
\item les éléments sont imbriqués (pas de superposition), 

\item présence d'\textbf{un seul élément racine},
\item etc  (se référer \`a la syntaxe XML).
\end{itemize}


Un document XML se doit d'être bien formé (sinon inexploitable).

\vfill
\info{
Pour tester si l'un de vos documents est bien formé :
un vérificateur syntaxique est disponible \`a l'adresse \url{http://www.w3schools.com/xml/xml_validator.asp}.\\
Vous pouvez également ouvrir votre fichier XML avec un navigateur web, 
intégrant généralement un parser 
XML (pour pouvoir \^etre parsé, un document doit \^etre bien formé).
}

 \end{frame}


%%%%%%%%%%%%%%%%%%%%%%%
\section{Formalisme XML: moins important (mais parfois utile)}



\begin{frame}[fragile]
\frametitle{Les entités et les références}

Les entités sont des symboles qui désignent des fragments. Utile pour réutiliser
du contenu, sans le répéter.

\vfill

Une entité est \alert{déclarée} (dans le prologue), puis \alert{référencée}.
 
\begin{lstlisting}[language=myXML]
<?xml version="1.0" encoding="ISO-8859-1"?>

<!DOCTYPE a [
  <!ENTITY  monNom "Charles Martel">
  <!ENTITY  maSignature SYSTEM "signature.xml">
]>

<a>
  Mon nom est &monNom;.

  &maSignature;
</a>
\end{lstlisting}

\end{frame}

%------------------------------------------------------------
\begin{frame}
\frametitle{Entités prédéfinies}

Tout un ensemble de symboles qui sont interprétées comme du marquage
 / balisage si on les utilise directement.

\vfill
 Si on les veut \alert{littéralement}, il faut utiliser des références d'entités.

\vfill
 
\begin{tabular}{l|l|l}
{\bf Déclaration} & {\bf Référence}&{\bf Symbole.}\\ 
\hline 
\texttt{<!ENTITY lt   "\&\#60;">}&\texttt{\&lt;}&\texttt{<}\\ 
\texttt{<!ENTITY gt   "\&\#62;">}&\texttt{\&gt;}&\texttt{>}\\ 
\texttt{<!ENTITY amp  "\&\#38;">}&\texttt{\&amp;}&\texttt{\&}\\ 
\texttt{<!ENTITY apos "\&\#39;">}&\texttt{\&apos;}&\texttt{'}\\ 
\texttt{<!ENTITY quot "\&\#34;">}&\texttt{\&quot;}&\texttt{"} \\ \hline
\end{tabular}

\end{frame}

%------------------------------------------------------------
\begin{frame}
\frametitle{Commentaires et instructions de traitement}

Les \alert{commentaires} peuvent être placés n'importe où dans la forme sérialisée.

\vfill

   \begin{center}
      \xmlinline{<!{-}- C'est un commentaire  {-}->}
    \end{center}

\vfill

  Ils apparaissent dans l'arbre DOM comme des n{\oe}uds de type \domtype{Comment}.

\begin{remark}
Très peu d'intérêt\,: on n'est pas censé \emph{lire} un document XML. Peut être remplacé par un
élément de type \domtype{element} et de nom \texttt{comment}.
\end{remark}
\vfill

\alert{Instructions de traitements}: commandes spécifiques à une application particulière.
\vfill
Typiquement, on demande à un processeur XSLT d'utiliser un programe \texttt{prog.xslt}\,:

\begin{center}
 \xmlinline{<?xml-stylesheet href="prog.xslt" type="text/xslt"?>}
\end{center}

\end{frame}


%------------------------------------------------------------
\begin{frame}[fragile]
\frametitle{Sections litérales}

Un analyseur (parseur) XML cherche à analyser tout le contenu pour détecter du marquage structurel.

\vfill 
\alert{Problème}: et si on ne veut \alert{pas} que le contenu soit analysé\,?

\begin{tinyverbatim}
<?xml version='1.0'?>
<program>
if ((i < 5) && (j > 6)) 
   printf("error");
</program>
\end{tinyverbatim}

\alert{Solution}: soit on utilise des entités (lourd), soit on insère la texte à protéger dans
une \alert{section littérale}.

\begin{tinyverbatim}
<?xml version='1.0'?>
<program>
<![CDATA[if ((i < 5) && (j > 6)) 
           printf("error");
]]>
</program>
\end{tinyverbatim}

\end{frame}

\begin{frame}
\frametitle{C'est tout pour l'instant\,!}

Ce qui précède suffit pour produire et comprendre à peu près tous les documents XML. 

\vfill

Autres aspects qui peuvent (quand même) s'avérer importants\,:

\begin{enumerate}
   \item \alert{Validation d'un document}. Comme déclarer qu'un document doit être conforme à une structure donnée? Comment le vérifier\,?
   
       $\Rightarrow$ typage de documents XML, avec les DTD ou XML Schéma (horreur!)
       
       \item \alert{Conflits de nom}. Comment bien interpréter une balise (nom d'élément) quand on assemble plusieurs documents\,?
       
       $\Rightarrow$  les \textit{namespaces} (espaces de noms).
         
     \item et autres aspects exotiques... 
   
\end{enumerate}

\end{frame}

\begin{frame}[fragile]
\frametitle{Ce qu'il faut retenir}

XML = tentative la plus complète de formaliser un modèle de représentation semi-structuré.

\begin{itemize}
  \item Très complet.
  \item Très lourd....  Syntaxe compliquée et redondante, éléments/attributs, beaucoup de détails peu utiles.
\end{itemize}
  
L'essentiel:

\begin{itemize}
  \item La structuration en arbre; l'existence d'une syntaxe pour coder cet arbre (la forme sérialisée);
  \item La notion de \alert{dialecte}, ou \alert{espace de nom}: des contraintes de nommage
     et de structuration pour lesquelles les documents ont une interprétation.
  \item JSON peut être vu comme une forme ultra-simplifiée; tend à remplacer XML 
         (sauf pour les documents dits "rédactionnels").
\end{itemize}

\end{frame}
